\subsection{Angular Accuracy at Driving Distances}
\label{sec:method_accuracy}
\todo{motivate why angular accuracy matters especially at driving distances}

We use a Pupil Labs Neon wearable eye tracker, which estimates gaze at 200\,Hz from two eye cameras without user calibration and records a front-facing scene camera (1600$\times$1200\,px, 30\,Hz, 103\dg$\times$77\dg{} field of view), eye-image streams and an IMU. All analyses use the uncorrected device gaze: wearer-specific offsets stored in the vendor app were removed before processing. Gaze samples are synchronised to scene frames and converted, like the target reference points, to azimuth/elevation directions using the camera intrinsics (positive right/up). The \emph{angular error} of a sample is the spherical angle between gaze and target direction; the \emph{bias} is the componentwise gaze-minus-target difference. Each participant is analysed separately and no samples are pooled across people. Participants are pseudonymised as P1--P6.
